\section{Related work}\label{sec:related}

\subsection{Underwater image enhancement}\label{sec:rw-enhancement}

Reviews divide the field into restoration, which inverts a model of image
formation, enhancement, which uses no model, and learned methods
\cite{schettini2010,wang2019review,lai2025,huang2026review}. Restoration
\cite{jaffe1990,akkaynak2018} corrects a loss that grows with range, which a
global operation cannot \cite{berman2020}, but the benchmarks that time
prior-based methods report seconds per frame
\cite{li2020uieb,peng2023,li2020uwcnn}. Model-free enhancement is the least
costly class: contrast-limited adaptive histogram equalization took 12.6~ms on a
300$\times$400 image against 0.43 to 4.58~s for eight prior-based and learned
methods \cite{liu2020ruie}. Nearest to the correction of
Section~\ref{sec:autovis} is \cite{bianco2015}, which shifts the chromatic
components by a grey-world estimate and stretches the luminance between
percentiles. The known failings of the class are an overcompensated red where
the scene has little red or few hues, amplified noise, and distorted colour
\cite{ancuti2012,ancuti2018,bianco2015,lai2025,wang2019review,liu2020ruie}; the
deadband and the limits of Section~\ref{sec:autovis} are aimed at the first two.
Learned methods account for most recent work \cite{cong2026}, and their
real-time claims are, as a rule, frame rates on a graphics processor at low
resolution \cite{islam2020,luo2026,zhu2026pic,yang2024lu2net}. Enhancement
applied frame by frame can flicker \cite{srinath2025,huang2026review}; the
remedies are a temporal filter on pixels \cite{ancuti2012}, networks that take
several frames \cite{du2024,xie2024}, optical-flow consistency losses in
training \cite{srinath2025} and an optimization over each finished video
\cite{zhu2025waterwave}, and none of them reports real-time operation on
embedded hardware. We found no evaluation of a first-order filter on the
parameters of a global correction, which is what Eq.~\eqref{eq:smooth} applies;
two reviews list temporal coherence at low cost as an open need
\cite{wang2019review,huang2026review}. The no-reference measures UIQM
\cite{panetta2016} and UCIQE \cite{yang2015} rise by construction with contrast
and saturation, and later studies find that they reward an overcorrected red and
disagree with observers \cite{li2020uieb,liu2020ruie,li2022limits,wang2019review};
we therefore evaluate on synthetic scenes whose undegraded form is known, a
practice that is criticized in its turn for not representing real water
\cite{berman2020}, and on the photographs of \cite{li2020uieb} against the
reference pictures that its observers chose. For object detectors, the unprocessed picture gave the best
accuracy for every detector tested \cite{wang2024detectors}, which is one reason
for \eeye{} to publish the camera's frame beside the pilot's; we found no study
of the effect of enhancement on piloting.

\subsection{Video delay and its transport}\label{sec:rw-latency}

In a pointing task a lag of 75~ms lengthened movement times measurably
\cite{mackenzie1993}, and in telesurgery a delay of 70~ms lengthened all four
manipulation tasks significantly \cite{nankaku2022}; for vehicles, reviews place
the onset of degradation between 170~ms and above 300~ms, depending on the task
\cite{chen2007,kamtam2024}. For underwater vehicles we found only impressions:
one experienced operator of a vehicle of the class considered here found that it
remained controllable up to about 100~ms of added delay and lost precision work
at 500~ms \cite{bauer2025}. For a small ROV observed through a network camera,
the delay from scene to display was 850 to 870~ms with H.264
\cite{kaknjo2018}, yet a chain instrumented stage by stage at 640$\times$480
spent 0.88~ms encoding, 0.02~ms on the network and 0.27~ms decoding of 19.67~ms
from scene to display \cite{bachhuber2018}. Pipelines left at their defaults
report 61 to 202~ms between encoder and decoder or from scene to display, each
without attribution to a stage
\cite{green2021,khalil2026,schimpe2022,jennehag2016,testouri2025}; where such a
delay was attributed, it was to a buffer, about 150~ms of 200~ms in the
receiver's jitter buffer \cite{baltaci2022}. The delays of the encoder are
documented: the zero-latency tuning of x264 reduced the delay through one 4K
chain from 41 to 82 frames to 7 \cite{arunruangsirilert2025}, and at a constant
bit rate the stream buffers of the encoder and the decoder are the largest delay
\cite{schreier2006,schreier2008}. Less is written about the receiver: a decoder
generally begins a frame only when all of it has arrived \cite{schreier2006},
and decoding several frames in parallel raises the frame rate without shortening
the delay of a frame \cite{chi2012} and in FFmpeg delays the output by a frame
for every thread beyond the first \cite{ffmpegmt}. In the work we read, that
delay is not measured, and the wait of a parser for the start of the next frame,
which Section~\ref{sec:transport} removes with a delimiter, is not reported. A
timestamp written into the pixels before encoding and read after decoding
\cite{green2021,khalil2026,baltaci2022}, the method of
Section~\ref{sec:verification}, resolves single frames but leaves out the camera
and the display, which optical methods include \cite{bachhuber2016,georg2020};
the figures of Table~\ref{tab:latency} are therefore the transport's share of
what the pilot sees.

\subsection{Vehicle software and the transport of frames}\label{sec:rw-software}

The low-cost vehicles built on the BlueROV2 share one architecture: an
autopilot, a Raspberry Pi companion computer and an Ethernet tether to a topside
computer \cite{vonbenzon2022,willners2021,edge2020,wilby2020,manzanilla2019}.
Those that add perception either add computers and are built on ROS or ROS~2
\cite{willners2021,edge2020,wilby2020,fernandez2026}, or send the video to the
surface and process it there
\cite{manzanilla2019,akram2022,akram2025,alinei2024}; none of them reports what
it costs to move a frame from one process to another. The frameworks used on
marine vehicles exchange messages
\cite{quigley2009,huang2010lcm,macenski2022,benjamin2010,martins2009,pinto2013},
and the cost of an image as a message has been measured outside the marine
field: the designers of ROS~2 state that its default transport is not effective
above 1~MB \cite{macenski2022}, and on a Raspberry Pi~4 a 640$\times$480 colour
image could not be passed between two processes at 20~Hz \cite{macenski2023}.
Transports through shared memory remove the copy through the kernel, each at a
price: generated message types \cite{wang2019tzc}, mutexes and condition
variables shared between processes \cite{iordache2021}, a kernel module
\cite{ishikawa2025}, or, in iceoryx, sizes of messages known in advance and
subscribers that write reference counts into the shared memory
\cite{pohnl2023}. The frame bus of Section~\ref{sec:bus} is a mapped file with a
documented layout and none of these dependencies, and pays instead a copy of
each frame on each side. Its reader follows the scheme of \cite{lamport1977}, in
which a writer never waits for readers, known as the sequence lock in Linux
\cite{boehm2012} and as the non-blocking write protocol in real-time systems
\cite{kopetz1993}; the bus is that protocol's form with several buffers, with
four, but has not its bound on the reader's repetitions, which rests on
guarantees of timing that Linux does not give. For the operators, NEON
implementations on ARM were 1.6 to 4.8 times as fast as the same code without
them, and the compiler's own vectorization left table look-ups scalar
\cite{pulli2012,khadem2023}, which is a reason to compose consecutive point
operators into one table (Section~\ref{sec:operators}). None of these sources
requires the vector and threaded paths to reproduce the scalar output exactly,
and of 58 pairs of scalar and vector routines of one widely used library that
were checked symbolically, 10 differed \cite{collingbourne2011}.
